\section{Methods}

\subsection{Design, registration, and reporting}

We conducted a systematic overview of reviews (meta-review) in accordance with a protocol prospectively registered in the International Prospective Register of Systematic Reviews (PROSPERO). An editable summary of the prespecified protocol is provided in Supplementary Appendix S1. Reporting was guided by PRISMA 2020 and the overview-specific PRIOR statement.\cite{page2021prisma,gates2022prior}

% AUTHOR ACTION: Insert the public PROSPERO CRD number and official record URL before submission.

\subsection{Information sources and search strategy}

We searched Web of Science, PubMed, IEEE Xplore, and Google Scholar. Supplementary searches covered reference lists of included reviews and the preprint servers TechRxiv, arXiv, bioRxiv, engrXiv, and medRxiv. Database-specific strategies combined controlled vocabulary, where available, and free-text terms for digital twins and adjacent digital patient constructs, human health and healthcare, and review methodology. Searches were limited to English-language reports published from 2020 onward. Initial searches were conducted in October 2025, with an updated search specified before final manuscript submission. The complete conceptual search strategy is reproduced in Supplementary Appendix S1.

% AUTHOR ACTION: Report the exact last-search date and complete database-specific strategies before submission.

\subsection{Eligibility criteria}

Eligible reports were systematic reviews, meta-reviews, umbrella reviews, scoping reviews, or structured narrative reviews that explicitly examined at least one generic or patient-specific digital model applied to human health. The model had to represent human physiology, biomechanics, or another biological process and address at least one review domain: conceptualization, construction, validation, application, or methodological evaluation. Eligible models could operate from the subcellular level through whole-body or in silico population levels.

We excluded primary studies without a structured review methodology; editorials, perspectives, commentaries, and conference abstracts without sufficient methodological detail; non-English reports; animal-only or non-biological models; and reviews limited to hospitals, logistics, infrastructure, communication systems, or medical devices without a patient-level modeling component. Reviews focused exclusively on behavioral or educational simulations without a physiological, biomechanical, or biological patient representation were also excluded.

\subsection{Study selection and inconsistency resolution}

Records were managed and screened in CADIMA. At title and abstract screening, each record was assessed independently by two reviewers using three questions: whether the report was a review with structured methodology, whether it addressed digital twins or any conceptually related term identified within the broader digital patient spectrum (such as virtual patient, in silico patient, or synthetic patient), and whether it concerned human health or healthcare. Responses were recorded as Yes, Unclear, or No. A record progressed to full-text review when there was no agreed No on any core criterion.

CADIMA-flagged discrepancies were handled using the prespecified adjudication procedure reproduced in Supplementary Appendix S2. Discrepancies that could not change the inclusion decision were verified and finalized by the review coordinator acting as a third reviewer. Discrepancies that could change eligibility were discussed by the two reviewers, with each providing a brief rationale; unresolved uncertainty favored inclusion, and the coordinator made the final decision when consensus was not reached. Full-text reports were assessed against the complete eligibility criteria using the entire report, and exclusion reasons were recorded for the study-flow diagram. Search and selection results are reported using PRISMA 2020 and PRIOR conventions.\cite{page2021prisma,gates2022prior}

% AUTHOR ACTION: The available protocol and workbooks do not record how many reviewers independently assessed full texts. Confirm and report that number before submission.

\subsection{Data extraction}

Data were extracted in CADIMA using a standardized form mapped a priori to the six review questions. The form was piloted on a subset of included reviews and refined before full extraction; one reviewer completed each extraction and a second reviewer verified the entry for accuracy. Bibliographic and administrative data were collected for every review, while question-specific fields addressed conceptual framing (RQ1), construction and modeling (RQ2), validation and fidelity (RQ3), application domains and functional roles (RQ4), barriers and infrastructural gaps (RQ5), and future research and standardization priorities (RQ6). The complete field groups, extraction guidance, and mapping to the review questions are provided in Supplementary Table~\ref{tab:extraction-framework}. Because this was an overview of reviews, the form also recorded whether each review assessed methodological quality or risk of bias in its included primary studies, which tool or criteria were used, and whether those judgments informed the review's synthesis or interpretation.

\subsection{Methodological quality and risk of bias}

The prespecified appraisal framework treated methodological quality and risk of bias as related but distinct constructs. AMSTAR~2 was selected for applicable systematic reviews of healthcare evidence, and SANRA for structured narrative reviews.\cite{shea2017amstar2,baethge2019sanra} Review-level bias concerns were organized around the four ROBIS process domains: study eligibility criteria, identification and selection of studies, data collection and study appraisal, and synthesis and findings.\cite{whiting2016robis} The CADIMA appraisal further considered uncritical validation or fidelity claims, limited generalizability, implementation barriers, and incomplete reporting of funding or conflicts of interest. Each applicable concern was to be recorded as high concern, some concern or unclear, or low concern. Appraisal judgments were not intended as automatic exclusion criteria; they were intended to contextualize the strength, consistency, and limitations of the narrative synthesis.

% AUTHOR ACTION: Complete and independently verify the appraisal fields before reporting appraisal results; the current workbook contains missing and non-standardized entries.

\subsection{Synthesis and reproducible analysis}

Synthesis was narrative and descriptive; no quantitative pooling was planned because the included reviews differed in scope, terminology, model architecture, and reported outcomes. Prespecified comparisons considered generic or archetype-based versus patient-specific models, static versus dynamic or continuously coupled models, and educational versus clinical decision-support applications.

All transformations from the extraction workbook to manuscript tables and figures were scripted in Python. Predefined, case-insensitive lexical rules were applied only to prespecified extraction fields; common null and negation phrases were removed before classification, each theme was counted at most once per review, and reviews could contribute to multiple non-mutually-exclusive categories. When a review described multiple biological levels, the highest reported level was used only for the hierarchy visualization. Automated checks evaluated required fields, duplicate normalized titles, publication-year conflicts, missingness, output completeness, and study-flow arithmetic. Source data were read without modification, and derived tables, figures, plotting data, and quality-control reports were written separately. Analysis code, classification rules, and derived outputs are available in the public reproducibility repository described in the Data and code availability section.
